% !PRINTED 652
\weberpage{652}
%<<CONT
der Wurzeln (2). Jede cyklische Function dieser Grössen ist
eine rationale Zahl. Was wir zu beweisen haben, ist, dass sich
die $x$ rational durch Einheitswurzeln ausdrücken lassen. Adjungiren
wir dem Körper $\Omega$ noch die $m^{\mathrm{te}}$ Einheitswurzel $r$, so entsteht
ein Körper $R(x, r)$, der die beiden Körper $\Omega = R(x)$ und
$R(r) = \Omega_m$ als Theiler enthält.

\gesperrt{Wenn eine Zahl} $\omega = \Phi(x, r)$ des Körpers $R(x, r)$ die
sämmtlichen Substitutionen $(x_0, x_1), (x_0, x_2), \ldots, (x_0, x_{m-1})$ gestattet,
so ist sie im Körper $\Omega_m$ enthalten; denn dann ist
\[
m\,\omega = \Phi(x_0, r) + \Phi(x_1, r) + \cdots + \Phi(x_{m-1}, r),
\]
und darin sind die Coëfficienten der Potenzen von $r$ nicht nur
cyklische, sondern sogar symmetrische Functionen der Wurzeln
(2), also rationale Zahlen. Dies gilt selbst noch in dem
Falle, dass die Gleichung für $x$ durch Adjunction von $r$ reducirt
wird.

Dies wenden wir an auf die \gesperrt{Resolventen}
\weq{3}{(r, x_0) = x_0 + r x_1 + r^2 x_2 + \cdots + r^{m-1} x_{m-1},}
und allgemeiner, wenn $s$ einen beliebigen ganzzahligen Exponenten
 bedeutet:
\weq{4}{(r^s, x_0) = x_0 + r^s x_1 + r^{2s} x_2 + \cdots + r^{(m-1)s} x_{m-1}.}

Durch diese Resolventen lässt sich $x_0$ selbst wieder rational
ausdrücken, nämlich:
\weq{5}{m x_0 = \sum\limits_{0,\,m-1}^{s} (r^s, x_0),}
und wenn wir also nachweisen können, dass alle diese Resolventen
 $(r^s, x_0)$ Kreistheilungszahlen sind, so haben wir unser
Ziel erreicht. Wenn wir
\weq{6}{(r^s, x_{\varkappa}) = x_{\varkappa} + r^s x_{\varkappa+1} + r^{2s} x_{\varkappa+2} + \cdots + r^{(m-1)s} x_{\varkappa+m-1}}
setzen, so geht $(r^s, x_{\varkappa})$ aus $(r^s, x_0)$ durch die Substitution
$(x_0, x_{\varkappa})$ hervor, und es besteht die fundamentale Relation
\weq{7}{r^{\varkappa s} (r^s, x_{\varkappa}) = (r^s, x_0).}

Alle diese Resolventen sind Zahlen des Körpers $R(x, r)$;
wir betrachten vorzugsweise die beiden folgenden Verbindungen:
\weq{8}{(r^s, x_0)^m = \omega_s,}
\weq{9}{(r^s, x_0)^{-1} (r, x_0)^s = \alpha.}
